%*******************************************************
% Introduction generale
% Chapitre non numerote, place avant la partie CV.
%*******************************************************
\pdfbookmark[1]{Introduction générale}{IntroductionGenerale}
\phantomsection
\addcontentsline{toc}{chapter}{\tocEntry{Introduction générale}}

\chapter*{Introduction générale}
\label{ch:introduction}
\markboth{Introduction générale}{Introduction générale}

% ====================================================================================
% GABARIT. Les trois hdr de reference (Rabaste, Ovarlez, Chatelain) enchainent
% les memes huit mouvements, dans cet ordre. Les blocs ci-dessous les suivent.
% Aucun des trois n'ouvre par « ce memoire presente »: l'accroche vient d'abord,
% la premiere personne n'arrive qu'au mouvement 5.
%
% Longueur visee: 6 a 8 pages. A ecrire APRES les bilans de chapitre, qui auront
% deja dit le detail --- sinon l'introduction double leur contenu et gonfle.
% ====================================================================================

% ============ TODO texte (~1 page) --- MOUVEMENT 1: L'ACCROCHE =====================
% Une accroche qui ne parle pas de soi. Rabaste ouvre sur l'etymologie de RADAR
% en 1942 et un amateur de palindromes; Chatelain sur le deluge de donnees, les
% puces a ADN et Efron; Ovarlez sur sa filiation intellectuelle.
%
% Proposition pour ce memoire --- un objet qui a change de statut:
%  - la matrice de covariance arrive comme statistique descriptive du XIXe
%    siecle (Pearson, puis Wishart en 1928);
%  - elle devient le parametre de nuisance d'un test de detection, ce qui
%    deplace la question de « la decrire » vers « l'estimer assez bien pour
%    decider »;
%  - elle finit comme une couche a l'interieur d'un reseau de neurones, qu'il
%    faut differentier --- donc comme un objet que l'on apprend, et non plus
%    seulement que l'on estime.
% Cet arc est celui du memoire et permet d'atteindre le \cref{ch:spdnet} des le
% premier paragraphe plutot qu'a la centieme page.
% ====================================================================================

% ============ TODO texte (~1,5 page) --- MOUVEMENT 2: POURQUOI LE PROBLEME EXISTE ==
% Le modele a etudier est la chaine causale d'Ovarlez: physique, puis
% consequence, puis besoin --- sans parler de soi. Chez lui: resolution plus
% fine -> moins de diffuseurs par case, variant d'une case a l'autre ->
% heterogeneite et impulsivite -> l'hypothese gaussienne tombe -> la regulation
% des fausses alarmes se degrade -> il devient *imperatif* d'avoir des outils
% robustes.
%
% La chaine de ce memoire est DOUBLE, et le dire tot est ce qui justifie la
% \cref{pt:context}:
%
%  BRANCHE 1 (robustesse). Resolution -> peu de diffuseurs par cellule ->
%  fouillis non gaussien, heterogene -> la scm et le cfar gaussien echouent
%  ensemble -> M-estimateurs, Tyler, modeles gaussiens composes
%  (\cref{ch:context-distributions}).
%
%  BRANCHE 2 (dimension). Nombre de variables croissant -> observations rares
%  relativement a la dimension ($\Nobs < \Nvar$) -> estimee singuliere ou mal
%  conditionnee (\cref{sec:context-definition-estimation}) -> structure et
%  regularisation -> et la structure est une contrainte GEOMETRIQUE, pas
%  seulement algebrique (\cref{ch:context-riemann}).
%
% Insister sur le fait que les deux branches ne sont pas independantes: c'est
% la these du memoire, annoncee ici et demontree au \cref{ch:spdnet}.
% ====================================================================================

% ============ TODO texte (~5 lignes) --- MOUVEMENT 3: LA PHRASE PIVOT ==============
% UNE phrase nommant le probleme unique dont tout le dossier est une facette.
% Rabaste: « peuvent etre consideres comme des facettes differentes d'un seul et
% meme probleme general: la detection de cibles faibles en environnement
% complexe ». Un jury cite cette phrase.
%
% Registre vise ici: estimer, contraindre et apprendre une matrice de covariance
% lorsque les observations sont rares, non gaussiennes, et vivent dans un espace
% courbe.
% Trois adjectifs -> trois chapitres d'etat des lieux -> trois contributions.
% A ciseler: c'est la phrase la plus relue du memoire.
% ====================================================================================

% ============ TODO texte (~0,5 page) --- MOUVEMENT 4: LE VOCABULAIRE ===============
% Rabaste forge « cibles faibles », dit immediatement que ce n'est PAS l'usage
% radar standard, et en donne trois sens en liste --- qui sont ses trois
% chapitres.
%
% Ici, nommer les trois tensions, puis dire qu'elles ne sont pas independantes:
%   (i)   $\Nobs$ contre $\Nvar$;
%   (ii)  gaussien contre elliptique / gaussien compose;
%   (iii) euclidien contre riemannien.
% Puis la formule qui porte le memoire, en une ligne: choisir une moyenne, c'est
% choisir une geometrie, donc un modele statistique. L'annoncer ici comme une
% these; \cref{prop:spdnet-moyennes-frechet} l'etablit cent pages plus loin.
% ====================================================================================

% ============ TODO texte (~2 pages) --- MOUVEMENT 5: UN PARAGRAPHE PAR CHAPITRE ====
% Un paragraphe par chapitre de contribution. Chacun doit dire QUATRE choses:
% le probleme, l'approche retenue, CE QUI LA DEMARQUE DE LA LITTERATURE, et LA
% THESE OU LE STAGE QU'ELLE A PRODUIT. Rabaste nomme cinq doctorants dans sa
% seule introduction.
%
%  - \cref{ch:robust-detection}: tests d'hypotheses robustes sur trois capteurs
%    (croix de Mills, sonar frontal multi-pings, series temporelles RSO), avec
%    le compromis selectivite/robustesse comme motif recurrent.
%    Nommer les collaborations sonar et le volet detection de changements.
%
%  - \cref{ch:spdnet}: trois endroits ou une moyenne euclidienne doit ceder la
%    place a une moyenne de Frechet --- la donnee, le batch, les poids.
%    Nommer Jafuno (these), Pautrel (M2 puis post-doc), et Gallet EN PRECISANT
%    qu'il s'agit d'une collaboration entre pairs et non d'un encadrement.
%
%  - \cref{ch:perspectives}: deux directions ou un objet de la
%    \cref{pt:context} n'a pas ete pousse jusqu'au bout.
%
% C'est ici que le dossier montre l'encadrement comme partie de la science, et
% pas seulement comme une rubrique du \cref{ch:cv}.
% ====================================================================================

% ============ TODO texte (~0,5 page) --- MOUVEMENT 6: LE CREDO =====================
% Une phrase disant quel chercheur on est. Ovarlez, explicitement: « je me suis
% toujours interesse aux lacunes potentielles des techniques couramment
% utilisees et souvent considerees A TORT comme valides ».
%
% Deux elements disponibles ici, et tous deux verifiables:
%
%  (a) LE MOTIF RECURRENT, a annoncer comme un fil que le lecteur recroisera:
%      l'hypothese gaussienne n'est jamais une simple approximation, c'est
%      toujours une geometrie deguisee.
%
%  (b) LA REPRODUCTIBILITE COMME POSITION DE METHODE, et non comme note
%      technique. Chaque figure de simulation porte son commit, sa graine et
%      tous ses parametres, et se rejoue depuis la ligne imprimee sous elle.
%      Dire POURQUOI: un resultat qu'on ne peut pas rejouer n'est pas un
%      resultat --- et, en ecrivant ce memoire, c'est precisement cet appareil
%      qui a fait apparaitre plusieurs erreurs de signe et de conjugaison dans
%      du code deja publie. Aucun des memoires de reference ne peut en dire
%      autant: c'est un element differenciant, il faut l'assumer comme tel.
% ====================================================================================

% ============ TODO texte (~1 page) --- MOUVEMENTS 7 ET 8: PLAN ET PREMIERE PERSONNE
% (7) Annoncer la structure, Y COMPRIS ce qu'elle a d'inhabituel.
%
%  - LE PARAGRAPHE QUI JUSTIFIE LA \cref{pt:context}. C'est le paragraphe le
%    plus important de l'introduction, parce qu'aucun des trois memoires de
%    reference n'a a le faire: aucun n'a pris le risque d'une partie
%    pedagogique autonome. Deux lectorats, chacun depourvu de ce que l'autre
%    possede: le traiteur de signal radar/sonar qui connait les modeles
%    elliptiques et la M-estimation mais pas le vocabulaire riemannien, et le
%    geometre ou praticien de l'apprentissage sur varietes dans la situation
%    inverse. Dire clairement que cette partie n'est pas un passage oblige mais
%    un point de rencontre, qu'un lecteur peut accepter un outil sur parole et
%    y revenir, et que la carte des dependances existe pour que personne n'en
%    lise plus que necessaire.
%    C'est ce paragraphe qui transforme une structure de 70 pages inhabituelle
%    en choix editorial assume.
%
%  - Les quatre parties, et pourquoi le \cref{ch:cv} est place la ou il est.
%
%  - Un renvoi d'une ligne vers \cref{ch:commentlire} pour les parcours de
%    lecture, afin que les deux chapitres se partagent le travail au lieu de se
%    recouvrir. NE PAS refaire ici la carte des dependances.
%
% (8) Le passage a la premiere personne, annonce explicitement. Rabaste le fait
% et cela se lit bien.
% ====================================================================================

% ====================================================================================
% TODO decision: cette introduction rend caduque toute mise en contexte
% scientifique presente dans \cref{ch:commentlire}. Une fois ce chapitre ecrit,
% relire CommentLire.tex et en retirer ce qui releve du « quelle est la
% question »: il ne doit plus y rester que du « comment circuler ».
% ====================================================================================
